
Large language models generate confident but factually incorrect statements, a phenomenon termed hallucination that undermines deployment in high-stakes domains. Detecting hallucinations before they reach users requires uncertainty quantification, yet current methods impose computational overhead. Semantic entropy clusters multiple sampled outputs for linguistic invariance but requires 5-10 forward passes per query. Probe-based methods train classifiers on hidden states, adding inference latency. Single-pass approaches using final-layer output entropy ignore internal processing dynamics that may distinguish factual retrieval from fabrication.

This work hypothesizes that the trajectory of representations across late layers carries discriminative signal. Under the Cross-Layer Trajectory Instability (CLTI) framework, factual retrieval follows stable attractor dynamics---representations converge monotonically toward a knowledge-grounded answer. Fabrication, lacking grounded knowledge, requires iterative cross-layer constraint satisfaction, producing higher trajectory instability.

Three single-pass metrics computed from layers 24--31 (8 late layers) of LLaMA-2-7B on TruthfulQA MC1 operationalize this hypothesis:
\begin{itemize}
\item \textbf{Normalized Trajectory Instability (NTI)}: variance of per-layer entropy normalized by mean entropy
\item \textbf{Convergence Monotonicity Index (CMI)}: proportion of layer transitions where answer similarity increases
\item \textbf{Representational Competition Index (RCI)}: top-token flip patterns across layers
\end{itemize}

Experiments validate the core claim while refuting secondary mechanisms. NTI achieves AUROC 0.5657 on 5-fold cross-validation, exceeding the 0.55 threshold specified a priori. Combined with CMI and final-layer entropy, detection improves by +7.1\% AUROC (Fisher combined p = 1.15e-05) over an intercept-only baseline. However, trajectory metrics fail on low-entropy predictions (AUROC 0.5136, 95\% CI includes chance), and RCI flip patterns appear in >90\% of both correct and incorrect responses.

Contributions:
\begin{enumerate}
\item Existence evidence that single-pass trajectory features discriminate hallucinations beyond output entropy
\item Mechanism refinement showing the signal is driven by high-uncertainty cases; confident hallucinations remain undetected
\item Negative results demonstrating RCI flip patterns are universal in transformer processing
\end{enumerate}

